\chapter{Methodology}
\label{ch:method}

\section{Threat model}
\label{sec:threat}

The attacker can insert passages into the document collection but cannot
modify existing passages, the retriever, the model or the prompt. This is the
standard corpus-poisoning assumption \cite{zou2025poisonedrag,zhong2023poisoning}.
For most attacks the attacker also knows the query they target and which
embedder the system uses, which is realistic for the small set of open
embedders in common use. The attacker does \emph{not} need query access to the
deployed system. The one exception is the optimised blocker baseline
(\cref{sec:res-blocker}), which needs black-box access by design.

The attacker's goal is \emph{availability}: the system should fail to give a
useful answer to a legitimate query that it would otherwise have answered. A
refusal, an empty or degenerate response, and a deflection all count as
denials. A wrong answer does not count as a denial, which is why answer
correctness is measured separately (\cref{sec:correctness}).

\section{System under test}
\label{sec:system}

\Cref{tab:config} gives the fixed pipeline. Generation is greedy, so every run
is deterministic given its inputs. Run-to-run variation is therefore confined
to the retrieval study, where it was measured (\cref{sec:threats}).

\begin{table}[ht]
\centering
\small
\caption{Fixed RAG configuration used in every run.}
\label{tab:config}
\begin{tabular}{@{}ll@{}}
\toprule
Component & Setting \\
\midrule
Embedder  & \texttt{all-MiniLM-L6-v2}, 384 dimensions \cite{reimers2019sbert} \\
Index     & FAISS HNSW, inner product; $M=16$, efConstruction 200, efSearch 50 \\
Retrieval & top-$k = 5$ \\
Decoding  & Greedy; 256 new tokens (2{,}048 for R1-Distill) \\
Precision & 4-bit quantised weights, all eight models \cite{dettmers2023qlora} \\
\bottomrule
\end{tabular}
\end{table}

\subsection{Matched model pairs}

Each pair shares pretrained weights and differs only in post-training
(\cref{tab:pairs}). Three pairs are true base-versus-aligned contrasts. The
fourth, \pair{llama-r1-8b}, contrasts two post-trained models: Llama-3.1-8B
Instruct and DeepSeek-R1-Distill-Llama-8B, which is reasoning-distilled from
the same base \cite{deepseek2025r1}. It measures a different thing, so it is
reported separately throughout and never pooled with the other three.

\begin{table}[ht]
\centering
\small
\caption{Model pairs. Alignment level is an ordinal label: 0 base,
1--2 instruction-tuned, 4 reasoning-distilled.}
\label{tab:pairs}
\begin{tabular}{@{}llll@{}}
\toprule
Pair & ``Base'' side & ``Aligned'' side & Levels \\
\midrule
\pair{llama-3.1-8b} & Llama 3.1 8B & Llama 3.1 8B Instruct & $0 \to 2$ \\
\pair{qwen-2.5-7b}  & Qwen 2.5 7B  & Qwen 2.5 7B Instruct  & $0 \to 2$ \\
\pair{mistral-7b}   & Mistral 7B v0.3 & Mistral 7B Instruct v0.3 & $0 \to 1$ \\
\pair{llama-r1-8b}  & Llama 3.1 8B Instruct & DeepSeek-R1-Distill-Llama-8B & $2 \to 4$ \\
\bottomrule
\end{tabular}
\end{table}

\subsection{Prompts}
\label{sec:prompts}

Aligned models are prompted through their chat template, and base models with
a plain completion prompt that includes two answered examples. The chat prompt
also contains a \emph{refusal licence}: ``if the context doesn't contain the
answer, say so briefly''. This default (\texttt{auto}) is what a deployer would
naturally do, but it confounds training with wording. \Cref{sec:res-prompt}
measures the confound directly with a $2 \times 2$ prompt-by-training ablation.
Every \texttt{auto}-prompt contrast in this report is therefore an \emph{upper
bound} on the alignment effect, and the prompt-matched figure is quoted
alongside it.

\section{Datasets}
\label{sec:datasets}

\begin{table}[ht]
\centering
\small
\caption{Corpora. HotpotQA was run and then withdrawn (\cref{sec:threats}).}
\label{tab:datasets}
\begin{tabular}{@{}lrrrl@{}}
\toprule
Corpus & Passages & Queries & Clean recall@5 & Status \\
\midrule
BEIR NQ \cite{kwiatkowski2019nq,thakur2021beir} & 501{,}000 & 1{,}000 & 0.760 & primary \\
BEIR FiQA-2018 \cite{maia2018fiqa}               & 57{,}638  & 648     & 0.409 & replication \\
\sout{BEIR HotpotQA} \cite{yang2018hotpotqa}     & 500{,}995 & 995     & ---   & withdrawn \\
\bottomrule
\end{tabular}
\end{table}

The NQ corpus is the 1{,}000 query-relevant (gold) passages plus 500{,}000
filler passages sampled at random from the full 2.68M-passage BEIR corpus.
They are sampled rather than taken as a prefix, because BEIR corpora are not
randomly ordered and a prefix would bias the embedding neighbourhood that the
retrieval-layer attacks operate in. FiQA is used whole. FiQA is financial
opinion QA and not drawn from Wikipedia, so it tests whether results transfer
across domain, corpus size and retrieval quality.

Before a corpus was used, a screening script checked whether its multiple gold
labels per query were \emph{redundant}, meaning each passage answers alone, or
\emph{complementary}, meaning all are needed. NQ and FiQA are redundant.
HotpotQA is complementary by construction, which is why it was withdrawn
(\cref{sec:threats}).

\section{Attack taxonomy}
\label{sec:attacks}

\Cref{tab:attacks} lists the 13 attacks. Full construction details are in
\cref{app:attacks}. The families target different stages of the pipeline. A
(semantic jamming) attacks the model's judgement through content. B (resource
exhaustion) tries to exhaust context or generation. C (algorithmic complexity)
manipulates retrieval geometry. D (chain-of-thought deadlocks) constructs
evidence that cannot be reasoned to a conclusion. None of the 13
uses gradients or query access to the target system. Passages are templated,
and the C family selects them by embedding proximity using the public
embedder.

\begin{table}[ht]
\centering
\small
\caption[Attack taxonomy]{Attack taxonomy. Docs = adversarial passages
injected per query.}
\label{tab:attacks}
\begin{tabularx}{\textwidth}{@{}llrX@{}}
\toprule
ID & Name & Docs & Idea \\
\midrule
\multicolumn{4}{@{}l}{\emph{A~--- Semantic jamming}} \\
A1 & Guardrail Triggering        & 5   & On-topic passage with a safety-trigger snippet appended (weapons, drugs, hacking) \\
A2 & Contradiction Flooding      & 8   & Formal-register passages, three contradicting claim pairs each \\
A3 & Authority Spoofing          & 3   & Passages marked \textsc{classified}, \textsc{confidential} or \textsc{restricted} \\
\multicolumn{4}{@{}l}{\emph{B~--- Resource exhaustion}} \\
B1 & Context Window Saturation   & 5   & $\sim$1{,}200-word padded passages that flood the context \\
B2 & Generation Loop Induction   & 5   & Stacked recursive, self-referential blocks \\
B3 & Multi-Retrieval Amplification & 10 & Passages carrying follow-up query hooks \\
\multicolumn{4}{@{}l}{\emph{C~--- Algorithmic complexity (retrieval layer)}} \\
C1 & Embedding Space Clustering  & 200 & Templated passages packed within a tight radius of the query embedding \\
C2 & Index Pollution             & 200 & Near-duplicate copies ($\times$5) with small embedding perturbations \\
C3 & Adversarial Embedding Perturbation & 10 & Zero-width and homoglyph token perturbations \\
\multicolumn{4}{@{}l}{\emph{D~--- Chain-of-thought deadlocks}} \\
D1 & Logical Contradiction Traps & 5   & Formally contradictory deductive arguments \\
D2 & Circular Reference Chains   & 6   & Three-passage citation loops, A$\to$B$\to$C$\to$A \\
D3 & Epistemic Uncertainty Amplification & 5 & Stacked contestedness and hedge markers \\
D4 & Infinite Qualification Traps & 5  & The answer behind a six-deep chain of prerequisites \\
\midrule
RAND & Random injection (control) & 200 & C1's budget without the clustering \\
\bottomrule
\end{tabularx}
\end{table}

\section{Outcome measures}
\label{sec:metrics}

Each run poses every query twice, once with the clean collection and once with
the attack's passages added. Each response is classified as a denial or an
answer (\cref{sec:classifier}).

\paragraph{Clean refusal floor.} The fraction of queries denied with no attack
present. It is a property of the model and prompt, not of the attack.

\paragraph{Attributable attack success rate.} Let $Q_{\text{ans}}$ be the queries
a side answered on the clean run. Then
\begin{equation}
  \text{ASR} \;=\; \frac{\lvert\{\,q \in Q_{\text{ans}} : q \text{ denied under attack}\,\}\rvert}{\lvert Q_{\text{ans}}\rvert}.
  \label{eq:asr}
\end{equation}
Conditioning on clean-answerable queries removes the refusal floor from the
numerator, so ASR counts only denials the attack caused. The cost is that the
denominator differs between sides and is self-selected, so it is always
reported (\cref{sec:threats}).

\paragraph{Risk difference.} $\text{RD} = \text{ASR}_{\text{aligned}} -
\text{ASR}_{\text{base}}$. Positive values mean the aligned model is more
vulnerable.

\section{Statistical protocol}
\label{sec:stats}

A run is classified by two tests, and the verdict needs both.
\begin{enumerate}
  \item \textbf{Fisher's exact test}, aligned against base, on the attacked denial
    counts. This asks whether the two models differ.
  \item \textbf{McNemar's exact test} \cite{mcnemar1947} on the aligned side,
    pairing each query's clean and attacked outcome. Let $c$ be queries the
    attack broke and $b$ queries it fixed. This asks whether the
    \emph{attack} is responsible, as opposed to a pre-existing floor.
\end{enumerate}
Fisher $p$-values are corrected with Benjamini--Hochberg
\cite{benjamini1995fdr} across all runs on a corpus. NQ and FiQA are separate
pools. The verdicts are:

\begin{description}[style=nextline,leftmargin=3.2cm]
  \item[Genuine paradox] Fisher significant after FDR, aligned more
    vulnerable, \emph{and} McNemar significant with $c > b$.
  \item[Floor artefact] Fisher significant, McNemar not. The models differ, but
    the attack is not what separates them.
  \item[Protective] Fisher significant, aligned \emph{less} vulnerable.
  \item[Not significant] Fisher not significant.
\end{description}

The floor-artefact category exists because an aligned model with a 35\%
refusal floor will out-refuse its base model under \emph{any} condition. That
is not evidence about the attack.

\section{Measurement instrument: the refusal classifier}
\label{sec:classifier}

Every denial number in this thesis comes from a rule-based classifier that
labels each response as an answer or one of five refusal types (epistemic,
explicit safety, unspecified, generation failure, and so on). Because it is the
instrument, it was validated before any result was reported.
\begin{itemize}
  \item 300 responses were hand-labelled, stratified across corpus, pair, side
    and condition, and deliberately weighted towards hard boundary cases.
  \item A dev/holdout split kept the rows used to diagnose classifier errors
    apart from those used to score it.
  \item A peer independent of the project labelled 50 rows from the written
    codebook without seeing the classifier's output.
\end{itemize}

\begin{table}[ht]
\centering
\small
\caption{Refusal classifier validation. Cohen's $\kappa$ \cite{cohen1960kappa}.}
\label{tab:classifier}
\begin{tabular}{@{}lrrr@{}}
\toprule
Comparison & $n$ & Agreement & $\kappa$ \\
\midrule
Classifier vs.\ primary labeller, holdout & 150 & 94.7\% & 0.884 \\
Classifier vs.\ independent annotator     & 50  & 88.0\% & 0.725 \\
Human vs.\ human (primary vs.\ independent) & 50 & 86.0\% & 0.674 \\
Human test--retest (primary labeller)     & 50  & 90.0\% & 0.774 \\
\bottomrule
\end{tabular}
\end{table}

The classifier agrees with the independent annotator better than the two humans
agree with each other. Reweighted to the population frequency of each response
class, its error rate over the whole corpus is 2.05\%. Its residual errors are
mostly false negatives (9.9\% against 1.1\% false positives), so it
\emph{under}-counts refusals and the measured alignment gap is, if anything,
conservative.

The 0.884 figure carries some circularity. The same person wrote the
classifier's patterns, produced the primary labels and adjudicated disputes. The
0.725 against the independent annotator is therefore the figure to quote where
one conservative number is needed. \Cref{app:classifier} lists every
disagreement with the independent annotator.

\section{Answer correctness}
\label{sec:correctness}

BEIR queries carry no short answers, so denial alone cannot say whether an
answer was \emph{right}. To check this, 656 of the 1{,}000 NQ queries were
joined, on exact normalised question text, to their NQ-open short answers
\cite{lee2019latent}. Each response is then scored correct (a gold alias
appears), wrong, or refused. Two matchers are reported. The strict matcher
requires an exact alias token sequence, and the lenient one requires $\geq$50\%
of an alias's content tokens. Every conclusion in \cref{sec:res-correctness}
holds under both.
